% migrate-to: lob-deep-learning
\subsection{Feature Engineering for Microstructure Alpha}

\begin{taxonomy} \hlt{Relative-Value Feature Segmentation}
\begin{enumerate}[label=\roman*.]
\setlength{\itemsep}{0pt}
\item \hlt{basis}: raw spread, carry-adjusted spread, z-score, percentile, and distance to arbitrage bounds;
\item \hlt{hedge quality}: rolling beta, cointegration residual, factor residual, and hedge-ratio instability;
\item \hlt{convergence clock}: time to delivery, expiry, fixing, rebalance, event date, or funding reset;
\item \hlt{carry}: funding rate, borrow fee, dividend, storage, roll yield, and collateral return;
\item \hlt{execution synchrony}: leg liquidity, queue priority, fill correlation, and latency between venues;
\item \hlt{stress capacity}: margin usage, drawdown to liquidation, short recall risk, and borrow availability.
\end{enumerate}
These features separate cheapness from tradability.
\end{taxonomy}

\begin{model} \hlt{Factor-Residual Arbitrage View}\\
For instruments $i$ and $j$, write $r_i=\beta_i^\top f+\epsilon_i$ and $r_j=\beta_j^\top f+\epsilon_j$. A relative-value signal should explain whether the spread move came from common factors $f$, hedge-ratio error $\beta_i-\beta_j$, or residual shocks $\epsilon_i-\epsilon_j$.
\end{model}

\begin{principle} \hlt{Alpha Needs Comparative Advantage}\\
A feature is useful only if it captures information or execution capacity that competitors do not fully exploit. A high in-sample score without an economic reason is usually a selection result, not an edge.
\end{principle}

\begin{taxonomy} \hlt{Microstructure Alpha Evaluation Segmentation}
\begin{enumerate}[label=\roman*.]
\setlength{\itemsep}{0pt}
\item \hlt{forecast}: relation to future midprice, spread, fill, or volatility;
\item \hlt{tradability}: spread, queue loss, borrow, fees, and venue access;
\item \hlt{capacity}: decay of performance with size and participation;
\item \hlt{stability}: sensitivity to regime, tick size, latency, and participant mix;
\item \hlt{crowding}: whether many agents can observe and monetize the same state.
\end{enumerate}
The feature is only the first filter; the executable strategy is the object being evaluated.
\end{taxonomy}

\begin{principle} \hlt{Arbitrage Labels Need Path Information}\\
A spread that later converges is not automatically a good label.\\
The trade may still be infeasible if the interim drawdown breaches margin, borrow, inventory, or risk limits.\\
Label both terminal convergence and maximum adverse excursion.
\end{principle}

\begin{taxonomy} \hlt{Arbitrage Event Label Segmentation}
\begin{enumerate}[label=\roman*.]
\setlength{\itemsep}{0pt}
\item \hlt{entry}: spread crosses estimated arbitrage bound after costs;
\item \hlt{confirmation}: flow suggests uninformed pressure or slow common-factor adjustment;
\item \hlt{failure}: specific news, borrow loss, deal break, or hedge-ratio regime change;
\item \hlt{exit}: convergence, stop-out, forced liquidation, expiry, or capital reallocation.
\end{enumerate}
\end{taxonomy}
